\subsection{Источники ошибок оптического энкодера}

Полная погрешность оптического углового энкодера формируется из нескольких основных составляющих:

\begin{itemize}
\item \textbf{Ошибка нанесения шкалы} --- отклонение реального положения штрихов от расчетного, возникающее в процессе изготовления растра, недостатки технологии.
\item \textbf{Эксцентриситет шкалы} --- смещение геометрического центра растра относительно оси вращения; порождает ошибку в гармонике первого порядка
\item \textbf{Деформации шкалы} --- ошибки формы, повторяющиеся кратно за оборот (вторая, третья и высшие гармоники)
\item \textbf{Ошибки считывающей головки} --- нестабильность зазора между головкой и растром, параллакс, температурный дрейф фотоприемников
\end{itemize}

Суммарная ошибка после установки энкодера (ошибка установки) включает ошибку системы <<шкала-головка>>, основные слагаемые которой описаны выше, а также дополнительные составляющие, внесенные условиями монтажа, например, люфты и дребезги креплений, подшипников и т.д. Именно эта величина является практически значимой характеристикой для конкретного применения.

\subsubsection{Оценка влияния на измерения в приближении малых ошибок}

Считая все возможные отклонения конструкции от идеала малыми, представим расстояние от оси вращения до шкалы растра в виде частичной суммы ряда Фурье по полярному углу до $n$-го члена \cite{renishaw_encoder_accuracy}. Обозначим $r_0$ --- средний радиус диска (номинальный радиус), $a_k$ -- амплитуда $k$-й гармоники ошибки, $\Phi_k$ --- фазовый сдвиг гармоники:

\begin{equation*}
    R(\theta) = r_0 - \sum_{k=1}^{n}a_k\cos(k\theta-\Phi_k)
\end{equation*}

Угловая ошибка измерения напрямую связана с ошибками формы и установки. Считывающая головка фактически фиксирует смещение $s$ по дуге кодового диска, которое связано с фактическим радиусом и угловым смещением:

\begin{equation*}
    ds = R(\theta)\cdot d\theta
\end{equation*}

Вычисленное энкодером угловое смещение есть величина смещения по дуге нормированная на средний номинальный радиус диска:

\begin{equation*}
    d\theta' = \frac{ds}{r_0} = \frac{R(\theta)}{r_0}\cdot d\theta
\end{equation*}

Полное угловое смещение может быть вычислено путем интегрирования:
\begin{equation*}
    \Delta \theta' = \int_{\theta_1}^{\theta_2}{\frac{R(\theta)}{r_0}d\theta}
\end{equation*}

Ошибка углового измерения:
\begin{multline*}
    E_{12} = \Delta \theta' - \Delta \theta = \\=-\frac{1}{r_0}\int_{\theta_1}^{\theta_2}{\sum_{k=1}^{n}{a_k\cos(k\theta-\Phi_k)}d\theta} = \\
    =-\frac{1}{r_0}\sum_{k=1}^{n}\frac{a_k}{k}\cdot\left[\sin(k\theta_2-\Phi_k)-\sin(k\theta_1-\Phi_k)\right]
\end{multline*}

Таким образом, амплитуда углового измерения для $k$-й гармоники составляет:

\begin{equation*}
    E_k = \frac{a_k}{k\cdot r_0}
\end{equation*}

Это означает, что гармоники высокого порядка оказывают все меньшее относительное влияние на точность при фиксированной амплитуде деформации --- при равных амплитудах первая гармоника (эксцентриситет) вносит наибольший вклад.

Для численной оценки амплитуды угловой ошибки эксцентриситета в работах \cite{eccentricity_error}, \cite{heidenhain-modular-angle-encoders} предлагается использовать аналогичную формулу:

\begin{equation*}
    \Delta \varphi[''] \simeq \frac{412}{D[\text{мм}]}\cdot e[\text{мкм}]
\end{equation*}

где $e$ --- величина сдвига (эксцентриситета) в микрометрах, $D = 2r_0$ --- номинальный диаметр диска в миллиметрах. Коэффициент в числителе есть не что иное как перевод размерностей:
$$
\frac{1~\text{рад}}{1''} \times \frac{1~\text{мкм}}{1~\text{мм}}\times\frac{D}{r_0}  \simeq 206265 \times 10^{-3}\times2 \simeq 412
$$

\subsection{Методы калибровки поворотных энкодеров}

Традиционно выделяют два класса методов калибровки: сравнительные и безэталонные --- т.е. методы \textit{самокалибровки}.

\subsubsection{Сравнительная калибровка}

В методах сравнительной калибровки измерения калибруемого энкодера сопоставляются с показаниями внешнего эталона более высокой точности. Роль таких эталонов традиционно выполняют полигональные зеркала с фотоэлектрическими автоколлиматорами \cite{24gon}, \cite{polygon}, калиброванные поворотные столы \cite{table}, \cite{table_guide} или лазерные гониометры \cite{gonio}.

Основные ограничения метода:
\begin{enumerate}
    \item Число точек калибровки ограничено числом граней полигонального зеркала или позиций индексного стола.
    \item Требует точного механического сопряжения калибруемого устройства с эталоном.
    \item Не может выполняться после монтажа энкодера в систему --- карта погрешностей неприменима к смонтированному энкодеру, поскольку не учитывает установочные погрешности конкретного монтажа.
\end{enumerate}

\subsubsection{Методы самокалибровки} \label{subsubsec:selfcal}

Методы самокалибровки не требуют внешнего эталона в процессе работы и могут применяться непосредственно в условиях эксплуатации. Подробный обзор и анализ алгоритмов самокалибровки
в оптических энкодерах представлен в работе \cite{kiryanov-overview}, на которую буду ссылаться далее.

\paragraph{Метод путевого усреднения (МПУ)}\label{par:avg_path_method} \cite{mpu-ionak}, \cite{mpu-bachish}, \cite{cramer} относится к методам самокалибровки с регулярной локализацией считывающих головок и в современной трактовке рассматривается как частный случай метода усреднения при равном делении (EDA, Equal Division Average).

Идея МПУ состоит в том, что $N$ считывающих головок размещаются равномерно по окружности, результирующий сигнал формируется усреднением их показаний:
$$
\bar{x}(\varphi)
=
\frac{1}{N}
\sum_{k=0}^{N-1}
x\!\left(\varphi + \frac{2\pi k}{N}\right).
$$

При таком усреднении в выходном сигнале подавляются все гармоники инструментальной погрешности, кроме кратных $N$.

Например, в документации \cite{heidenhain-modular-angle-encoders} Heidenhain описывает вычисление усредненной позиции по двум  головкам для модульных энкодеров ERA/ERP. В документе явно указано, что при этом «устраняются все нечетные гармоники погрешности нарезки, а также погрешности от эксцентриситета установки» \footnote{\dots. The accuracy improvement actually attained by this in practice strongly
depends on the installation situation and the application. All eccentricity errors (reproducible errors due to mounting errors,
non-reproducible errors due to radial eccentricity of the bearing) and all uneven-numbered harmonics of the graduation error are eliminated. \cite{heidenhain-modular-angle-encoders}}. Это объясняется тем, что диаметрально противоположное размещение двух головок является минимальной конфигурацией для устранения эксцентриситета установки. При $N = 2$ из ряда Фурье исчезают все нечетные гармоники, включая доминирующую первую. Это следует из того, что при смещении аргумента на $\pi$ нечетные составляющие меняют знак и при сложении взаимоуничтожаются. 

\paragraph{Метод усреднения при равном делении (EDA)} \cite{watanabe-eda} является базовым и наиболее строго обоснованным методом самокалибровки для многоголовочных энкодеров с равномерным размещением считывающих головок.

Самокалиброванное значение угловой позиции можно записать в виде
$$
Q_i^{C}
=
\frac{1}{N_S}
\sum_{j=1}^{N_S}
Q_{i,j},
$$

где $Q_{i,j}$~--- показание $j$-й головки в $i$-й угловой позиции, а $N_S$~--- число считывающих головок.
Для получения калибровочной кривой отдельной головки $j$ сравнивают ее сигнал с усредненным сигналом ансамбля, т.~е. рассматривают разность
$$
c_j(\varphi)
=
x_j(\varphi) - \bar{x}(\varphi).
$$

В спектральном представлении это означает, что в калибровочной кривой отсутствуют гармоники, кратные $N_S$, тогда как амплитуды остальных гармоник воспроизводятся практически без искажения.

Моделирование, проведенное в работе \cite{kiryanov-overview}, показывает, что ансамбли считывающих головок с регулярной локализацией обеспечивают существенно более точное восстановление калибровочных кривых по сравнению с методами, основанными на нерегулярном размещении головок.

\paragraph{Метод Масуды} \cite{masuda} был предложен как компромисс между точностью и аппаратной сложностью многоголовочного датчика.

Вместо большого числа равномерно распределенных головок используется ограниченное число считывающих головок, расположенных в специально выбранных угловых позициях, а итоговая калибровочная кривая формируется по последовательности каскадных усреднений.

Удобно интерпретировать метод как рекурсивную схему вида
$$
S_{r+1}(\varphi)
=
\frac{1}{2}
\Bigl(
S_r(\varphi) + S_r(\varphi + \Delta_r)
\Bigr),
$$
где на каждом шаге $r$ используется свой угловой интервал $\Delta_r$ (в классических схемах анализируются шаги $90^\circ$, $45^\circ$, $22.5^\circ$, $11.25^\circ$ и т.~д.).

На промежуточных стадиях такая обработка приводит к последовательному исключению гармоник с номерами $4$, $8$, $16$, $32$ и т.~д., что формально подтверждает способность метода подавлять высокие гармоники при относительно малом числе головок.

Моделирование в работе \cite{kiryanov-overview}, показывает, что при всех достоинствах по подавлению характеристических гармоник метод Масуды существенно искажает амплитуды остальных гармоник.

\paragraph{Мультикомбинационный метод Ватанабэ} \cite{watanabe-selfa}
развивает идеологию EDA, но использует несколько независимых наборов считывающих головок с различным числом элементов в каждом наборе.

Пусть $r$-й набор содержит $N_r$ головок, равномерно размещенных по окружности; его частный усредненный сигнал можно записать как
$$
y_{N_r}(\varphi)
=
\frac{1}{N_r}
\sum_{k=0}^{N_r-1}
x\left(\varphi + \frac{2\pi k}{N_r}\right),
$$
а общий сигнал ансамбля~--- как
$$
y(\varphi)
=
\frac{1}{R}
\sum_{r=1}^{R} y_{N_r}(\varphi),
$$
где $R$~--- число независимых наборов.
При выборе форматов наборов, например $2$–$3$–$5$ или $6$–$7$ и т.~п., достигается подавление гармоник, кратных наименьшему общему кратному соответствующих чисел (для формата $2$–$3$–$5$ ожидаемо подавляются гармоники, кратные $30$).
Такой подход позволяет с относительно небольшим суммарным числом головок покрыть больший диапазон низких гармоник спектра погрешности.

В то же время моделирование, проведенное в работе \cite{kiryanov-overview} для формата $2$–$3$–$5$, показывает, что амплитуды доступных гармоник воспроизводятся с разной точностью: гармоники, присутствующие во всех наборах, детектируются более точно, тогда как гармоники, возникающие только в части наборов, имеют заметную систематическую ошибку.

\paragraph{VEDA-метод (Virtual Equal Division Averaged)} \cite{ishii-veda}
представляет собой попытку воспроизвести эффект EDA без установки большого числа \textit{физических} считывающих головок.
Вместо реального равномерного ансамбля используются несколько головок в специальных угловых позициях, а недостающие каналы восстанавливаются вычислительно, то есть виртуально.

Преимущество VEDA состоит в снижении аппаратной сложности: часть информации, которую в EDA дает фактическое увеличение числа головок, здесь появляется за счет более сложной обработки и накопления данных.

В работе \cite{kiryanov-overview} все методы с нерегулярной локализацией головок (включая VEDA) демонстрируют худшие метрологические характеристики по сравнению с методами на основе регулярной локализации (МПУ и EDA): амплитуды гармоник калибровочных кривых воспроизводятся менее точно, а остаточный вклад инструментальной погрешности растра остается выше.
Таким образом, VEDA-метод целесообразно рассматривать как компромисс, обеспечивающий экономию по числу головок ценой усложнения алгоритмов и увеличения погрешности.

\subsection{Запретные зоны телескопов и модель наведения}
\label{sec:forbidden}

\subsubsection{Определение и классификация запретных зон}
\label{subsec:forbidden_def}

Запретной зоной телескопа называется область угловых положений на небесной сфере (область значений координат монтировки), в которой телескоп не может находиться в рабочем положении по физическим, конструктивным или программным причинам.

Совокупность запретных зон определяет \textit{рабочее поле} инструмента и является характеристикой наблюдательного комплекса~\cite{ESO_Horizon}. Принято различать три основных класса запретных зон.

\paragraph{Механические запретные зоны.}Порождаются конечными угловыми пределами осей монтировки: упорами, укладкой кабелей, корпусом купола, а также возможными столкновениями трубы с элементами конструкции.

Для альт-азимутальной монтировки основная механическая запретная зона~--- зенитная слепая зона (\textit{zenith blind spot}): двигатели, приводящие телескоп в движение, имеют ограничение по скорости вращения телескопа по азимуту, что означает, что в зените будет область, в которой двигатели не смогут вращать телескоп достаточно быстро (или достаточно точно), чтобы продолжать отслеживать объект. Это приводит к появлению <<слепой зоны>>, которая обычно находится в пределах одного градуса от зенита~\cite{TPoint_Wallace}. Кроме того, ограничение снизу (минимальная высота, как правило $15.0^\circ$-$20.0^\circ$) задается требованием к оптической толщине атмосферы и механическими упорами монтировки. 

\paragraph{Астрофизические и атмосферные ограничения.} Включают максимальную допустимую зенитную дистанцию, ограничение по яркости Луны и Солнца, горизонтные маски рельефа и горизонта. Они не связаны с механикой напрямую, однако вводятся в систему управления как программные запреты.

\paragraph{Зависимые ограничения двухосных монтировок.} Класс, характерный для немецких экваториальных монтировок (GEM). Вследствие кинематики GEM достижимость точки с данными координатами зависит от того, на каком меридиане находится труба, поскольку одна и та же точка неба может быть достигнута в двух ориентациях трубы (\textit{East side} и \textit{West side}), разделенных меридианным переворотом~\cite{ATUS_Stuttgart}\cite{EQMOD_Limits}.

\subsubsection{Запретные зоны в немецких экваториальных монтировках.} Немецкая экваториальная монтировка оснащена двумя ортогональными осями: полярной осью и осью склонений. Труба закрепляется с одной стороны оси склонений, уравновешиваясь противовесами с другой. Из-за асимметрии конструкции допустимый диапазон угла по полярной оси ограничивается не только механическими упорами ($t \in [-t_{\max}, +t_{\max}]$), но и условием отсутствия физического контакта трубы и навесного оборудования с опорной колонной монтировки~\cite{TPoint_Wallace}.

Принципиальный запрет состоит в следующем: при наблюдении объекта с $\delta < \delta_{\text{pier}}$, когда труба пересекает меридиан, монтировка обязана выполнить \textit{меридианный переворот} (\textit{meridian flip})~--- переброс трубы на противоположную сторону с поворотом оси склонений на $180^\circ$. В окрестности меридиана существует \textit{запретная зона часового угла} шириной $\Delta t_{\text{flip}}$, в которой монтировка не может корректно завершить слежение без смены стороны; для робота-телескопа или автоматического планировщика наблюдений эта зона должна быть явно задана и контролироваться системой управления~\cite{EQMOD_Limits}\cite{EKOS_Flip}. Существенно, что ограничение по часовому углу и склонению зависимы~\cite{TPoint_Wallace}.

\paragraph{Пример: КГО ГАИШ МГУ.}

Кавказская горная обсерватория ГАИШ МГУ (плато Шатджатмаз, высота 2100~м, $43^\circ44'$~с.~ш., $42^\circ40'$~в.~д.) располагает двумя телескопами с принципиально различными типами монтировок~\cite{Shatsky2020}.

\textbf{2.5-метровый телескоп SAI2.5} оснащен \textit{альт-азимутальной} монтировкой с прямыми приводами (direct drive) Etel и контроллером движения DeltaTau PMAC2. Для данной монтировки характерны два основных класса запретных зон: 

\begin{enumerate}
    \item Ограничение по высоте, программно $h_{\min} \approx 15^\circ$--$20^\circ$; 
    \item Зенитная слепая зона
\end{enumerate}

Радиус зоны вокруг зенита на широте $\varphi$, в пределах которой достигается допустимая максимальная скорость привода по азимуту, определяется соотношением:

$$
r_z \approx \arctan\left(\frac{\dot{A}_{\max}}{\omega \cos\varphi}\right),
$$

где $\omega= 15.04''/{s}$~--- угловая скорость суточного вращения, $\dot{A}_{\max}$~--- максимальная угловая скорость по азимуту, обеспечиваемая приводом. Слежение через зенит требует специального алгоритма переброса~\cite{TPoint_Wallace}.

\textbf{60-см робот-телескоп RC600} оснащен \textit{экваториальной} монтировкой DDM160 (ASA Astrosysteme GmbH) с прямыми приводами~\cite{Shatsky2020}. Для данного инструмента запретные зоны имеют структуру, характерную для немецкой монтировки:

\begin{enumerate}
    \item Меридианный переворот при пересечении меридиана;
    \item Программные ограничения по часовому углу с учетом механических упоров и требований к безопасности оборудования
\end{enumerate}

При автоматическом режиме работы, в котором RC600 функционирует с мая 2019~г., точный учет запретных зон особенно важен: планировщик наблюдений обязан заблаговременно прогнозировать момент меридианного переворота и корректно выбирать момент его выполнения.

\subsubsection{Влияние погрешностей энкодера на запретные зоны}
\label{subsec:encoder_forbidden}

Точность энкодера на оси телескопа непосредственно влияет на ширину программных буферов безопасности вокруг запретных зон. Пусть суммарная погрешность измерения углового положения оси составляет $\sigma_{\text{enc}}$. Тогда минимальный буфер, гарантирующий, что монтировка не войдет в запретную зону вследствие ошибки измерения, должен удовлетворять условию
\[
\sigma_{\text{buf}} \geq k \cdot \sigma_{\text{enc}},
\]
где $k$~--- коэффициент безопасности. При наличии незамеченных низкочастотных систематических погрешностей (эксцентриситет установки, биение диска) $\sigma_{\text{enc}}$ может достигать десятков угловых секунд, что вынуждает увеличивать буфер и тем самым сокращать доступную область неба. Таким образом самокалибровка энкодера, описываемая в \ref{subsubsec:selfcal}, имеет дополнительное прикладное значение: снижение $\sigma_{\text{enc}}$ позволяет уменьшить буферы безопасности и расширить рабочее поле инструмента.

\subsection{Модель наведения телескопа по звездам}
\label{subsec:pointing}

\subsubsection{Структура погрешностей наведения и их источники}
\label{subsubsec:pointing_errors}

Точность наведения телескопа определяется не только качеством угловых энкодеров, но и рядом других систематических погрешностей, которые имеют принципиально различную физическую природу~\cite{TPoint_Wallace}\cite{EngineersView}. Здесь и далее $t$ обозначает часовой угол, $\delta$ --- склонение, а $Z$ --- зенитное расстояние. Принято выделять следующие основные классы погрешностей согласно ~\cite{TPoint_Wallace}:

\paragraph{Индексные ошибки.} Нулевые поправки осей, термины IH и ID в TPoint, представляют собой постоянные смещения нулевых положений энкодеров относительно истинного нуля координатной системы монтировки. Они возникают, например, вследствие поворота шкалы энкодера при монтаже, теплового расширения элементов крепления, неточности начальной синхронизации. В отличие от случайных погрешностей, индексные ошибки дают одинаковый вклад во всех угловых положениях и могут быть устранены наведением на опорную звезду.

\paragraph{Геометрические ошибки неперпендикулярности осей.} Обусловлены реальными механическими отклонениями конструкции. В частности, неперпендикулярность оси склонений к полярной оси (терм NP) приводит к угловым ошибкам, пропорциональным $\tan\delta$:

$$
\Delta t_{\text{NP}} = \mathtt{NP} \cdot \tan\delta.
$$

Аналогично неперпендикулярность оптической оси трубы к оси склонений (терм CH) создает ошибку наведения, которую невозможно устранить простой коррекцией нуля:

$$
\Delta t_{\text{CH}} = \mathtt{CH} \cdot \sec\delta.
$$

Эти ошибки не связаны с энкодером: они возникают в механической структуре монтировки и при фиксированном состоянии инструмента остаются стабильными.
    
\paragraph{Ошибки ориентации полярной оси.} Термы ME и MA отражают реальное отклонение оси вращения монтировки от направления на полюс мира. В модели TPOINT вертикальная составляющая этого отклонения описывается термом ME:

$$
\Delta t_{\text{ME}} = \mathtt{ME}\,\sin t \,\tan\delta,
\qquad
\Delta \delta_{\text{ME}} = \mathtt{ME}\,\cos t ,
$$

а горизонтальная (лево--право, по азимуту установки полярной оси) --- термом MA:

$$
\Delta t_{\text{MA}} = \mathtt{MA}\,\cos t \,\tan\delta,
\qquad
\Delta \delta_{\text{MA}} = -\,\mathtt{MA}\,\sin t .
$$

Эти ошибки вносят систематический вклад на всем небе и в идеальном случае должны устраняться механической юстировкой монтировки. Хотя модель наведения способна оценить и частично компенсировать их программно, предпочтительным является уменьшение коэффициентов ME и MA именно механической настройкой полярной оси, поскольку это снижает остаточные систематические ошибки и уменьшает риск нежелательного вращения поля~\cite{TPoint_Wallace}.

\paragraph{Гнутия трубы.} Терм TF описывает деформацию оптической трубы под действием силы тяжести по закону Гука:
    
$$
\Delta Z_{\text{TF}} = \mathtt{TF} \cdot \sin Z,
$$

Гнутие является типичным примером погрешности, \textit{физически независимой} от погрешности энкодера: оно влияет на положение изображения в фокальной плоскости, но не на показание угловых датчиков монтировки. Тем не менее в рамках модели наведения гнутие входит в ту же линейную задачу наименьших квадратов, что и ошибки энкодеров.

\subsubsection{Метод калибровки наведения по звездам (pointing model)}
\label{subsubsec:pointing_method}

Процедура построения модели, реализованная в пакете TPoint~\cite{TPoint_Wallace}, состоит в следующем.

Оператор проводит сеанс наблюдений эталонных звезд, равномерно распределенных по небосводу (рекомендуется 30--50 звезд~\cite{TPoint_Wallace}). Для каждой звезды в момент наблюдения регистрируется пара данных: \textit{истинное положение} звезды (прямое восхождение и склонение из каталога, приведенные к эпохе с учетом прецессии, нутации, аберрации и рефракции) и \textit{механические показания} угловых датчиков монтировки $(t_{\text{meas}}, \delta_{\text{meas}})$. Вектор невязки
$$
\boldsymbol{\varepsilon}_i = \bigl(t_{\text{meas},i} - t_{\text{true},i},\;\;
\delta_{\text{meas},i} - \delta_{\text{true},i}\bigr)
$$
для каждой звезды является суммой всех погрешностей наведения в данной точке неба. Модель наведения~--- это аналитическое представление функции $\boldsymbol{\varepsilon}(t, \delta)$ в виде линейной комбинации физически обоснованных термов:
\[
\boldsymbol{\varepsilon}(t,\delta) = \sum_{k} c_k \, \mathbf{T}_k(t, \delta),
\]
где $\mathbf{T}_k(t, \delta)$~--- базисные функции, соответствующие известным геометрическим и механическим погрешностям (IH, ID, NP, CH, ME, MA, TF и др.), а коэффициенты $c_k$ определяются методом наименьших квадратов. Стандартом является шестичленная геометрическая модель для экваториального телескопа \texttt{IH ID NP CH ME MA}; в нее дополнительно включаются гнутия, периодические ошибки редуктора и прочие эффекты по мере выявления систематических остатков на диагностических графиках.

В этом смысле модель наведения выполняет функцию глобальной калибровки всей системы телескопа~--- от истинного направления на объект до положения изображения в выбранной точке фокальной плоскости.

Важно отметить, что модель наведения не устраняет \textit{нестационарные} погрешности~--- тепловые дрейфы узлов, гистерезис редуктора, изменение гнутия при смене навесного оборудования. TPoint предоставляет команду \texttt{GHYST} для диагностики гистерезиса и \texttt{GAM} для мониторинга изменений ориентации полярной оси; обнаружение нестационарности требует повторной калибровки.

\subsection{Алгоритмы субпиксельной точности}

\subsubsection{Актуализация проблемы}

При регистрации изображения фотоприемной матрицей непрерывный оптический сигнал преобразуется в дискретную последовательность отсчетов. Каждый пиксель линейки усредняет световой поток, падающий на его площадку, и выдает одно числовое значение. Пространственный шаг такой выборки равен физическому размеру пикселя.

Согласно теореме Найквиста–Шеннона, для однозначного восстановления непрерывного сигнала по его дискретным отсчетам достаточно, чтобы частота дискретизации не менее чем вдвое превышала максимальную частоту в спектре сигнала. Иными словами, теорема устанавливает минимально необходимую частоту выборки, при которой сигнал может быть восстановлен без потерь. Однако из этого условия не следует, что при выполнении теоремы позиция особенности сигнала определяется лишь с точностью до одного пикселя. Если сигнал является полосоограниченным — то есть его спектр не содержит компонент выше некоторой пространственной частоты, — то соседние отсчеты связаны между собой аналитическим образом через функцию sinc-интерполяции. Это означает, что значение сигнала в любой промежуточной точке между двумя пикселями теоретически полностью определяется всей совокупностью имеющихся отсчетов.

На практике применение полной sinc-интерполяции по всей длине сигнала вычислительно избыточно и нецелесообразно. Вместо этого используется информация о форме профиля яркости вблизи локальной особенности (пика или фронта), которая достаточно полно характеризуется небольшим числом соседних отсчетов. Их совокупность несет информацию о субпиксельном положении особенности. Различные субпиксельные методы отличаются тем, каким образом они извлекают эту информацию \cite{blais_rioux_peak}.

Субпиксельная точность определения позиции ограничена уровнем шума в сигнале. При отношении сигнал/шум, типичном для ПЗС-матриц в условиях умеренной засветки, стандартные методы позволяют определять положение фронта или пика с СКО порядка десятых долей пикселя. Это соответствует улучшению пространственного разрешения почти на порядок по сравнению с пиксельной точностью без какого-либо аппаратного изменения системы.

\subsubsection{Метод центра масс \cite{blais_rioux_peak}}

Позиция пика определяется как взвешенное среднее координат пикселей в окрестности
максимума:

$$
x_{\rm peak} = x_0 + \frac{\displaystyle\sum_{i=-N}^{N} I(x_0 + i)\cdot i}
                           {\displaystyle\sum_{i=-N}^{N} I(x_0 + i)}
$$

Параметр $N$ должен быть достаточно мал для исключения влияния шума и соседних пиков.

\subsubsection{Параболическая интерполяция \cite{phase_matching_sampling} \cite{blais_rioux_peak}}

Трем соседним точкам вокруг максимума вписывается парабола $y = ax^2 + bx + c$, вершина которой дает субпиксельную позицию:

$$
x_{\rm max} = -\frac{b}{2a}
$$

Формула для расчета при равноотстоящих отсчетов (шаг 1 пикс.):

$$
x = x_0 + \frac{I(x_0 - 1) - I(x_0 + 1)}
               {2\,I(x_0 - 1) - 4\,I(x_0) + 2\,I(x_0 + 1)}
$$

Метод широко используется при детектировании фронтов в задачах обработки лазерных линий и кодовых растров.

Если парабола вписывается не в исходный сигнал, а в его производную, в таком случае вершина параболы соответствует точке перегиба оригинального сигнала~— то есть субпиксельной позиции фронта между битами.

\subsubsection{Метод Блэ–Риу \cite{blais_rioux_peak}}

\begin{enumerate}
\item Дифференцирование сигнала через асимметричный КИХ-фильтр 4-го порядка (BR2):
$$
H_1(z) = 1 + z^{-1} - z^{-3} - z^{-4}
$$
или 8-го порядка (BR4):
$$
H_2(z) = 1 + z^{-1} + z^{-2} + z^{-3} - z^{-5} - z^{-6} - z^{-7} - z^{-8}
$$
Оба фильтра выполняют одновременно сглаживание и дифференцирование.

\item Линейная интерполяция нулей функции в точке, где производная меняет знак между
двумя соседними отсчетами $A \geq 0$ и $B < 0$. 
Поправка к положению нуля относительно абсциссы точки $A$:
$$
\Delta = \frac{A}{A - B}
$$

\item Итоговая позиция пика (фильтр вносит задержку $(N-1)/2$).

$$\mathrm{pos} = x_A + \frac{A}{A-B} - \frac{N-1}{2}$$

\end{enumerate}

\subsubsection{Метод нулевого перехода второй производной (метод Стегера) \cite{steger_subpixel}}

\begin{enumerate}
\item Свертка с $\dfrac{d}{dx}G_\sigma(x)$~--- первой производной гауссова ядра
с шириной $\sigma$.

\item Нулевые переходы производной отклика $r'(x)$ при условии $r''(x) < 0$ указывают
на максимумы яркости (центры светлых полос).

\item Субпиксельная позиция нулевого перехода определяется линейной
интерполяцией двух соседних отсчетов с разными знаками:
$$
x_{\rm sub} = x_i + \frac{r'(x_i)}{r'(x_i) - r'(x_{i+1})}
$$
\end{enumerate}

Параметр $\sigma$ должен соответствовать ожидаемой ширине полосы растра.

\subsubsection{Субпиксельная детекция фронтов методом Кэнни–Деверне \cite{gioi_canny_devernay}}

\begin{enumerate}
\item Оператор Кэнни определяет позицию фронта как максимум нормы градиента вдоль
направления градиента.

\item Коррекция квадратичной интерполяцией нормы градиента в трех точках вдоль направления:
$$
\delta = -\frac{g_{+1} - g_{-1}}{2\,(g_{+1} - 2\,g_0 + g_{-1})}
$$
где $g_{-1},\, g_0,\, g_{+1}$~--- значения нормы градиента в соседних позициях.
\end{enumerate}

\subsubsection{Метод апсемплинга (передискретизации) \cite{phase_matching_sampling}}

Изображение перед детектированием увеличивается по числу пикселей в $M$ раз
с помощью интерполяции (билинейной, бикубической, sinc-интерполяции). После этого
стандартный пиксельный детектор работает на более мелкой сетке, обеспечивая
субпиксельную точность относительно исходного изображения. Число пикселей возрастает
как $M^2$, что существенно увеличивает вычислительную нагрузку. Поэтому апсемплинг
целесообразен только при нечастых вызовах или для верификации.

\subsection{Ограничения на кодовую последовательность псевдослучайного оптического растра}
\label{sec:pseudorand_code_constraints}

\subsubsection{Структура кодовой последовательности}\label{subsubsec:code_structure}

Кодовая дорожка оптического растра с псевдослучайным кодом имеет не имеет четкой структуры, однако последовательность бит образует подпоследовательность бинарной \textit{последовательности де Брейна}\footnote{Последовательностью де Брейна $B(k,~n)$ называют циклический порядок $a_1a_2\ldots a_t$ из конечного алфавита $\{0,\ldots, k-1\}$, такой, что любое слово длины $n$ из данного алфавита встречается в качестве ее подпоследовательности ровно один раз} длины $N$. 

Подпоследовательность в том смысле, что такая кодовая последовательность длины $N$ --- циклическая бинарная строка длины $1 \leq N \leq 2^n$, в которой каждое $n$-битное слово встречается не более одного раза. При $N = 2^n$ это полная последовательность де Брейна \cite{debruijn_cutdown}.

Для повышения надежности распознавания кодовых слов выполнялся перебор вариантов кодовых последовательностей при заданных значениях длины кодового слова $n$ и полной длины последовательности $N$. Критерием выбора служила минимизация максимальной длины однородного участка последовательности, состоящего из подряд идущих нулей или единиц.

Основными конструктивными параметрами кодового диска являются число бит кода $n$, ширина одного бита $w$ на внешнем радиусе, внутренний радиус кодовой дорожки $R$ и ее радиальная толщина $\Delta R$.

Положим геометрию конструкции фиксированной и сосредоточимся на ограничениях на кодовую последовательность для достижения необходимых характеристик

\subsubsection{Угловое разрешение}

Обозначим $\sigma(k)$ СКО определения положений бит, $k$ параметр среднего числа пикселей матрицы на один бит кода, $\lambda$ линейная плотность пикселей на единицу длины изображения растра.

Угловое разрешение выражается через эти величины

\begin{equation} \label{eq:resolution}
    \boxed{\rho = \dfrac{\sigma}{\lambda R}}
\end{equation}

\subsubsection{Число допустимых кодов}

Структура псевдослучайного кода такова, что каждый бит однозначно соответствует уникальному коду. В лучшем случае, если удалось правильно полностью разметить все поле зрения матрицы, то положение одного кода и даже одного его конкретного бита точно определяет ориентацию растра

В поле зрения матрицы должно попадать не менее $n+1$ бита, чтобы возможно было распознать $n$-битный код.

Поскольку геометрия фиксирована, пусть поле зрения матрицы определяется по длине $L \ll 2\pi R$ обозримого участка растра. Таким образом, пропорционально, полное число $N$ бит последовательности есть

$$N \geq \frac{2\pi (R+\Delta R)}{L}\cdot(n+1)$$

Полное число бит последовательности, являющейся подпоследовательностью де Брейна, также удовлетворяет соотношению 

$$2^{n-1} \leq N \leq 2^n$$

Таким образом:
\begin{equation}
\label{eq:code_len}
    \begin{cases} n-1\leq \log_2N\leq n,\\
    N \geq \frac{2\pi (R+\Delta R)}{L}\cdot(n+1)
    \end{cases} \Rightarrow \boxed{n \geq \log_2\left(\dfrac{2\pi (R+\Delta R)}{L}\cdot(n+1)\right)}
\end{equation}

\subsubsection{Геометрическое условие размещения растра на диске}
\label{subsec:geometry_track}

Геометрия кодовой дорожки на диске должна обеспечивать корректное считывание при любом угловом положении.

Рассмотрим случай, когда линейка ориентирована под малым углом $\alpha$ к касательной середины дорожки и граница поля зрения касается внешнего радиуса дорожки (см. рис \ref{fig:estimate_dR}).
\begin{figure}[H]
    \centering
    \includegraphics[width=0.65\linewidth]{images/оценка на ширину дороки.pdf}
    \caption{Рассматриваемый случай}
    \label{fig:estimate_dR}
\end{figure}

Чтобы весь отрезок $L$ целиком оставался в пределах кодовой дорожки должно выполняться условие

По теореме косинусов внутри треугольника $\Delta ABC$:

\begin{equation*}
(R+\Delta R)^2 = \left(R+\frac{\Delta R}{2}\right)^2 + \left(\frac{L}{2}\right)^2 -2\cdot\left(R+\frac{\Delta R}{2}\right)\cdot\frac{L}{2}\cdot \cos\left(\frac{\pi}{2}+\alpha\right)    
\end{equation*}

Считая, что $\Delta R \ll R$ получаем
\begin{equation}\label{eq:estimate_dR}
    \boxed{\Delta R \gtrsim L\cdot \dfrac{\sin \alpha + \frac{L}{4R}}{1 - \frac{L}{2R}\sin\alpha}}
\end{equation}

\subsubsection{Ограничения применимости и влияние аберраций линзы}
\label{subsec:lens_imperfections}

\paragraph{Дисторсия и нелинейность увеличения.}
Реальная линза имеет переменное увеличение по полю зрения. В результате шаг пикселя, приведенный к плоскости растра, не одинаковпо всей длине активной зоны. Это вносит систематическую погрешность в соответствие между номером пикселя и угловым положением.

\paragraph{Аберрации и конечная глубина резкости.}
Дефокусировка и аберрации приводят к уширению изображения бита на матрице. Если эффективный размер пятна рассеяния превышает размеры пикселя, пространственное разрешение ограничивается оптикой, а не матрицей

\paragraph{Непараллельность оси линзы и плоскости растра.}
Наклон оптической оси относительно плоскости диска порождает трапециевидную дисторсию: масштаб изображения различен у ближнего и дальнего края поля зрения.

\paragraph{Практический вывод.}
Все перечисленные выше условия являются \textit{необходимыми}, но в присутствии перечисленных аберраций не являются достаточными для гарантированного считывания кода. На этапе первичного прототипирования рекомендуется вводить дополнительный запас по ширине бита и полю зрения сверх теоретического минимума, окончательное определение характеристик проводится экспериментально.
